% Generated draft: requires mathematical and source review.
\documentclass[11pt]{article}
\usepackage[margin=25mm]{geometry}
\usepackage{fontspec}
\setmainfont{Latin Modern Roman}
\usepackage{amsmath,amssymb,mathtools,unicode-math}
\setmathfont{Latin Modern Math}
\usepackage{enumitem,needspace,xurl,hyperref}
\setlength{\parindent}{0pt}
\setlength{\parskip}{5pt}
\setlength{\emergencystretch}{3em}
\newcommand{\N}{\mathbb{N}}
\newcommand{\Z}{\mathbb{Z}}
\newcommand{\Q}{\mathbb{Q}}
\newcommand{\R}{\mathbb{R}}
\newcommand{\C}{\mathbb{C}}
\newcommand{\F}{\mathbb{F}}
\newcommand{\A}{\mathbb{A}}
\newcommand{\PP}{\mathbb{P}}
\newcommand{\E}{\mathbb{E}}
\newcommand{\eps}{\varepsilon}
\newcommand{\dd}{\,\mathrm{d}}
\begin{document}
\section*{UnsolvedMath ID: 30004100}
\textbf{Source problem code:} catalogue entry
\subsection{Definitions and mathematical statement}
Let \(R\) denote the finite number of moments retained by a reconstruction algorithm. The record also calls \(R\) the truncation order; the correspondence between order and number of retained moments is not specified. The moments are quantities used to reconstruct an underlying mathematical object, but the record does not identify that object, its domain, the moment functions, or the reconstruction algorithm. In particular, it does not specify whether the moments are known exactly or estimated from observations.

The supplied research question is: how should \(R\) be selected in practical reconstruction computations so that retaining too many moments does not lead to overfitting or inaccurate reconstruction?

The supplied literature assessment describes the intended finite-sample objective as selecting the truncation order while controlling approximation error, estimation error, and overfitting. However, no sampling model, admissible class of objects, loss function, definition of overfitting, error threshold, or probabilistic guarantee is supplied. Consequently, this record specifies a research objective but does not determine a unique quantified mathematical assertion or a precise performance requirement for a selection rule. No universal guarantee is asserted here.

The source registry labels the problem open. Its supplied dated assessment says that convergence as the number of moments grows is understood, while a generally valid finite-sample selection rule controlling the stated errors and overfitting is unavailable. These are supplied curation claims, not independently verified mathematical or literature conclusions.
\subsection{Short English statement}
Choose how many moments \(R\) a practical reconstruction algorithm should retain, with the aim of avoiding overfitting and inaccurate reconstruction when too many moments are used. A precise guarantee requires specifications absent from the record, including the reconstructed object, moments, observation model, reconstruction method, and error criterion.
\subsection{Sources}
\begin{itemize}
\item Supplied local UnsolvedMath registry record: permanent ID \(30004100\), title ``Adaptive Moment Selection Without Reconstruction Overfitting,'' exact source problem code \texttt{catalogue entry}, in the Oberwolfach Reports Open Problems collection. The supplied proposed year is 2019; no individual proposer is identified.
\item Supplied source citation: \emph{Uncertainty Quantification (2020)}, Oberwolfach Reports, DOI 10.4171/owr/2019/12, pp.\ 695--772. Source URL and supplied literature-source link, labelled ``Oberwolfach report'': \url{https://doi.org/10.4171/owr/2019/12}. The background attributes the collection to workshop reports and their open-problem discussions. The supplied citation has not been independently consulted or verified for this draft.
\item Literature assessment checked at \texttt{2026-08-22}: the supplied status is ``open.'' The supplied assessment states that convergence as the number of moments grows is understood, but that no generally valid finite-sample rule is available for selecting the truncation order while controlling approximation error, estimation error, and overfitting. This is supplied dated curation, not an independently verified assessment; the record explicitly requests independent verification of time-sensitive status and cited claims.
\item The registry supplies the statement label \texttt{corrected\_verified} and says that mathematical content was preserved while notation, typography, duplicate clauses, or omitted definitions were repaired. This is a supplied editorial qualification, not independent mathematical verification of the present draft.
\end{itemize}
\end{document}
